\section{The Most Holy Eucharist: Real Presence, Sacrifice, and Banquet}
\label{sec:eucharist}

The Most Holy Eucharist is the summit of Christian Initiation and the living center of the Church. Whereas Baptism brings the believer to new birth and Confirmation strengthens him with the seal of the Holy Spirit, the Eucharist contains the Author of grace Himself---Body, Blood, soul, and divinity---under the appearances of bread and wine. St.~Thomas Aquinas states its singular preeminence directly: ``Absolutely speaking, the sacrament of the Eucharist is the greatest of all the sacraments \dots\ First of all because it contains Christ Himself substantially'' (\emph{ST}~III, q.~65, a.~3). At the altar, the Church enters into the one redeeming sacrifice of Calvary made sacramentally present, receives the living Bread from heaven, and is united into one Mystical Body in charity until the Lord returns.

\subsection{Prefigurations, Promise, and Institution in Sacred Scripture}
\label{subsec:eucharist-scripture}

The Old Testament prepared for the Eucharistic mystery through five converging figures. When Abraham returned from battle, ``Melchisedech, the king of Salem, bringing forth bread and wine, for he was the priest of the most high God, blessed him'' (Gen~14:18--19). Psalm~109:4 swore that the Messiah would be ``a priest for ever according to the order of Melchisedech,'' a title the Letter to the Hebrews applies directly to Christ's eternal priesthood (Heb~5:6; 7:24--27). On the night of the Exodus, Israel slew an unblemished Passover lamb, marked the doorposts with its blood, ate its roasted flesh with unleavened bread, and kept that night ``for a memorial to you \dots\ with an everlasting observance'' (Exod~12:5--14, 27, 46). Through the desert, God rained manna from heaven (Exod~16:4, 15), which St.~Paul reads as ``spiritual food'' foreshadowing Christ (1~Cor~10:3--4). At Sinai, Moses sprinkled sacrificial blood upon the people, declaring, ``This is the blood of the covenant, which the Lord hath made with you'' (Exod~24:8). And Malachias foretold a universal offering replacing the temple sacrifices: ``from the rising of the sun even to the going down, my name is great among the Gentiles, and in every place there is sacrifice, and there is offered to my name a clean oblation'' (Mal~1:11).

These figures met their fulfillment in Jesus of Nazareth. After multiplying the loaves in Galilee, He delivered the Bread of Life discourse at Capharnaum (John~6:48--59 in the Clementine Vulgate and Douay-Rheims, corresponding to John~6:48--58 in Greek and modern editions, because the Vulgate divides v.~51 into vv.~51--52 and runs one verse higher thereafter). When His hearers stumbled at the realism of eating His flesh, Jesus only deepened His solemn insistence:
\begin{quotation}
I am the living bread which came down from heaven. If any man eat of this bread, he shall live for ever: and the bread that I will give is my flesh, for the life of the world. \dots\ Amen, amen, I say unto you: except you eat the flesh of the Son of man and drink his blood, you shall not have life in you. He that eateth my flesh and drinketh my blood hath everlasting life: and I will raise him up in the last day. For my flesh is meat indeed: and my blood is drink indeed. He that eateth my flesh and drinketh my blood abideth in me: and I in him. (John~6:51--52, 54--57 DR [6:51, 53--56 Gk])
\end{quotation}
He closed by lifting this Bread infinitely above the manna: ``Not as your fathers did eat manna and are dead. He that eateth this bread shall live for ever'' (John~6:59 DR [6:58 Gk]).

On the night before He suffered, at the Paschal supper, Jesus fulfilled that promise by instituting the Eucharistic sacrifice and the priesthood to perpetuate it. In the Synoptic accounts (Matt~26:26--28; Mark~14:22--24; Luke~22:19--20), He took bread, blessed, broke, and gave it to His disciples, saying, ``Take ye and eat. This is my body,'' and over the chalice declared, ``For this is my blood of the new testament, which shall be shed for many unto remission of sins'' (Matt~26:26--28), adding the sacrificial charge, ``Do this for a commemoration of me'' (Luke~22:19). St.~Paul records the same tradition received from the Lord:
\begin{quotation}
For I have received of the Lord that which also I delivered unto you, that the Lord Jesus, the same night in which he was betrayed, took bread, And giving thanks, broke and said: Take ye and eat: This is my body, which shall be delivered for you. This do for the commemoration of me. In like manner also the chalice, after he had supped, saying: This chalice is the new testament in my blood. This do ye, as often as you shall drink, for the commemoration of me. For as often as you shall eat this bread and drink the chalice, you shall shew the death of the Lord, until he come. (1~Cor~11:23--26)
\end{quotation}
From that Real Presence St.~Paul draws the stark warning that whoever eats or drinks unworthily becomes ``guilty of the body and of the blood of the Lord,'' eating and drinking judgment to himself by ``not discerning the body of the Lord'' (1~Cor~11:27--29). Earlier in the same letter, he binds communion directly to the Church's unity: ``The chalice of benediction which we bless, is it not the communion of the blood of Christ? And the bread which we break, is it not the partaking of the body of the Lord? For we, being many, are one bread, one body: all that partake of one bread'' (1~Cor~10:16--17). From Pentecost onward, the Jerusalem Church persevered ``in the doctrine of the apostles and in the communication of the breaking of bread and in prayers'' (Acts~2:42).

\subsection{The Patristic Confession: Realism, Sacrifice, and Unity}
\label{subsec:eucharist-fathers}

Across East and West, the Fathers confessed with one astonished voice that the consecrated gifts are truly the crucified and risen Flesh and Blood of Christ, offered as a pure sacrifice and received as the food of immortality.

On his way to martyrdom around the year 110, St.~Ignatius of Antioch urged the Ephesians to obey the bishop and presbytery, ``breaking one and the same bread, which is the medicine of immortality, and the antidote to prevent us from dying, but [which causes] that we should live for ever in Jesus Christ'' (\emph{Ephesians}~20, shorter recension). Confronting Docetists who denied Christ's real flesh, he points straight to the altar:
\begin{quotation}
They abstain from the Eucharist and from prayer, because they confess not the Eucharist to be the flesh of our Saviour Jesus Christ, which suffered for our sins, and which the Father, of His goodness, raised up again. (St.~Ignatius of Antioch, \emph{Smyrnaeans}~7, shorter recension)
\end{quotation}
In the chapter immediately following (\emph{Smyrnaeans}~8), Ignatius binds that Eucharistic realism directly to apostolic hierarchy, declaring that a valid and proper Eucharist is celebrated under the bishop or one to whom the bishop has entrusted it.

A generation later in Rome, St.~Justin Martyr described the Sunday liturgy---readings from the apostles and prophets, exhortation, common prayer, bread and wine mixed with water, thanksgiving answered by the people's ``Amen,'' and distribution by deacons to those present and absent (\emph{First Apology}~65, 67)---and grounded Eucharistic realism directly in the Incarnation:
\begin{quotation}
For not as common bread and common drink do we receive these; but in like manner as Jesus Christ our Saviour, having been made flesh by the Word of God, had both flesh and blood for our salvation, so likewise have we been taught that the food which is blessed by the prayer of His word, and from which our blood and flesh by transmutation are nourished, is the flesh and blood of that Jesus who was made flesh. (St.~Justin Martyr, \emph{First Apology}~66)
\end{quotation}

Against Gnostic disdain for matter, St.~Irenaeus of Lyons identified the Eucharist as the universal ``pure sacrifice'' of Mal~1:11 offered to the Creator out of His creation (\emph{Adversus haereses}~IV.18.1, 4), and linked it directly to bodily resurrection: ``For as the bread, which is produced from the earth, when it receives the invocation of God, is no longer common bread, but the Eucharist, consisting of two realities, earthly and heavenly; so also our bodies, when they receive the Eucharist, are no longer corruptible, having the hope of the resurrection to eternity'' (\emph{Adversus haereses}~IV.18.5; cf.~V.2.2--3).

In fourth-century Jerusalem, the fourth and fifth \emph{Mystagogical Catecheses} (ascribed traditionally to St.~Cyril of Jerusalem, with modern scholarship also weighing authorship by his successor John~II of Jerusalem) warn neophytes not to let bodily taste overrule Christ's word:
\begin{quotation}
Since then He Himself declared and said of the Bread, \emph{This is My Body}, who shall dare to doubt any longer? \dots\ He once in Cana of Galilee, turned the water into wine, akin to blood, and is it incredible that He should have turned wine into blood? \dots\ the seeming bread is not bread, though sensible to taste, but the Body of Christ; and that the seeming wine is not wine, though the taste will have it so, but the Blood of Christ. (\emph{Mystagogical Catechesis}~IV.1--2, 9)
\end{quotation}
In the fifth lecture, the catechist describes the invocation of the Holy Spirit upon the gifts, ``that He may make the Bread the Body of Christ, and the Wine the Blood of Christ; for whatsoever the Holy Ghost has touched, is surely sanctified and changed,'' and explains that over ``that sacrifice of propitiation'' the Church offers ``Christ sacrificed for our sins, propitiating our merciful God for them as well as for ourselves,'' for both the living and the dead (\emph{Mystagogical Catechesis}~V.7--8, 10).

In Milan, St.~Ambrose appeals to God's creative omnipotence in \emph{De mysteriis}~VIII--IX:
\begin{quotation}
Let us prove that this is not what nature made, but what the blessing consecrated, and the power of blessing is greater than that of nature, because by blessing nature itself is changed. \dots\ Shall not the word of Christ, which was able to make out of nothing that which was not, be able to change things which already are into what they were not? \dots\ Before the blessing of the heavenly words another nature is spoken of, after the consecration the Body is signified. He Himself speaks of His Blood. Before the consecration it has another name, after it is called Blood. And you say, Amen, that is, It is true. (St.~Ambrose, \emph{De mysteriis}~IX.50, 52, 54)
\end{quotation}

St.~John Chrysostom taught his flock to trust Christ's word above sight: ``For His word cannot deceive, but our senses are easily beguiled. \dots\ Since then the word saith, `This is my body,' let us both be persuaded and believe, and look at it with the eyes of the mind'' (\emph{Homilies on Matthew}, Hom.~LXXXII.4). Christ Himself remains the true Consecrator: ``He that then did these things at that supper, this same now also works them. We occupy the place of servants. He who sanctifieth and changeth them is the same'' (\emph{Homily}~LXXXII.5). Preaching on 1~Cor~10:16--17, he identifies the chalice with Calvary---``This which is in the cup is that which flowed from His side, and of that do we partake''---and asks, ``For what is the bread? The Body of Christ. And what do they become who partake of it? The Body of Christ: not many bodies, but one body'' (\emph{Homilies on 1~Corinthians}, Hom.~XXIV.3--4).

St.~Augustine articulated the ecclesial sacrifice in \emph{De civitate Dei}~X.6: ``This is the sacrifice of Christians: we, being many, are one body in Christ. And this also is the sacrifice which the Church continually celebrates in the sacrament of the altar, known to the faithful, in which she teaches that she herself is offered in the offering she makes to God.'' In Tractate~26.11--13 on John, he exclaims---``O mystery of piety! O sign of unity! O bond of charity!''---while distinguishing bare outward eating from spiritual communion: ``the sacrament is one thing, the virtue of the sacrament another \dots\ he that eateth within, not without; who eateth in his heart, not who presses with his teeth.''

Finally, St.~John of Damascus gathers the Greek synthesis in \emph{De fide orthodoxa}~IV.13. The glorified body does not spatially leave heaven; rather, ``the bread itself and the wine are changed into God's body and blood'' by the presence and invocation of the Holy Spirit. Rejecting mere symbolism outright, he writes: ``The bread and the wine are not merely figures of the body and blood of Christ (God forbid!) but the deified body of the Lord itself: for the Lord has said, `This is My body,' not, this is a figure of My body.''

\subsection{Transubstantiation, Concomitance, and the Mode of Presence}
\label{subsec:eucharist-transubstantiation}

To protect the literal truth of Christ's words ``This is my body'' and ``This is my blood,'' the Church formulated a precise metaphysical vocabulary. The Fourth Lateran Council in \emph{Firmiter credimus} (1215, DS~802) confessed that Christ's Body and Blood are truly contained under the species of bread and wine (\latin{in sacramento altaris sub speciebus panis et vini veraciter continentur}), ``the bread being transubstantiated into the body, and the wine into the blood, by divine power'' (\latin{transsubstantiatis pane in corpus, et vino in sanguinem potestate divina}). The Council of Trent in Session~XIII (1551, chs.~1--5, cans.~1--11) defined that Jesus Christ is contained \emph{truly, really, and substantially} (\latin{vere, realiter ac substantialiter}), excluding any reduction to bare sign, figure, or power, and declared that the Catholic Church most aptly names the conversion of the whole substance of bread into the Body and the whole substance of wine into the Blood, the species alone remaining, \textbf{transubstantiation} (\latin{transsubstantiatio}; Sess.~XIII, ch.~4, can.~2).

St.~Thomas Aquinas (\emph{ST}~III, qq.~75--77), the \emph{Roman Catechism} (Part~II), Pope Paul~VI (\emph{Mysterium Fidei}~38--49), St.~John Paul~II (\emph{Ecclesia de Eucharistia}~15), Pope Benedict~XVI (\emph{Sacramentum Caritatis}~11), and the \emph{Catechism of the Catholic Church} (CCC~1373--1377, 1390) preserve five strict distinctions.

\paragraph{1. Transubstantiation is a unique substantial conversion.}
In natural changes, some underlying matter endures as a new form replaces the old. In the Eucharist, no underlying matter of bread or wine remains, nor is their substance merely annihilated so Christ's Body can move into an empty space (\emph{ST}~III, q.~75, aa.~2--3). By God's creative omnipotence, the \emph{entire} substance of bread passes into the substance of Christ's Body, and the \emph{entire} substance of wine into His Blood:
\begin{quotation}
The whole substance of the bread is changed into the whole substance of Christ's body, and the whole substance of the wine into the whole substance of Christ's blood. Hence this is not a formal, but a substantial conversion; nor is it a kind of natural movement: but, with a name of its own, it can be called ``transubstantiation.'' (\emph{ST}~III, q.~75, a.~4)
\end{quotation}
This excludes both \emph{consubstantiation} (bread remaining alongside Christ) and \emph{impanation}, and sets the Eucharist apart from the other six sacraments, which elevate human souls accidentally or establish a relational bond without altering human substance.

\paragraph{2. Transubstantiation is strictly distinct from the determining words of consecration.}
Christ's words of institution---``This is my body,'' ``This is the chalice of my blood''---are the \emph{instrumental sacramental form} acting within the Eucharistic Prayer (\emph{ST}~III, q.~78, aa.~1, 4). Transubstantiation itself is the \emph{divinely wrought substantial conversion} effected through that instrument, never a mere shift in communal meaning or liturgical purpose. In \emph{Mysterium Fidei} (1965, \S\S38--49), Pope Paul~VI distinguished Christ's Eucharistic presence \emph{par excellence} from His other real modes of presence, rejecting theories of ``transignification'' or ``transfinalization'' separated from ontological conversion. New meaning follows from new substantial reality; words without substantial conversion would leave bread as bread.

\paragraph{3. Transubstantiation is strictly distinct from any change of accidents.}
Every sensible property of bread and wine---the \latin{species}, including color, taste, texture, aroma, weight, chemical composition, and dimensive quantity---remains completely unaltered (\emph{ST}~III, q.~75, a.~5; q.~77, aa.~1--2; Trent, Sess.~XIII, ch.~4, can.~2). The senses are never deceived, because the real sensible accidents of bread and wine genuinely remain; what changes is the invisible substance attained ``by faith alone'' (\emph{ST}~III, q.~75, a.~1). Since bread's substance no longer exists and Christ's glorified Body cannot be the subject of bread's whiteness or roundness, divine power directly sustains the remaining dimensive quantity without a substantial subject, and all other sensible qualities inhere in that quantity (\emph{ST}~III, q.~77, aa.~1--2). Breaking or consuming the Host divides the species alone, leaving Christ's impassible Body wholly undivided.

\paragraph{4. Presence \latin{ex vi verborum} versus real concomitance.}
Under the consecrated bread, Christ's \textbf{Body} is present \latin{ex vi verborum} (``by the force of the words'' or \latin{ex vi sacramenti}), and under the consecrated wine His \textbf{Blood} is present \latin{ex vi verborum} (\emph{ST}~III, q.~76, aa.~1--2). Yet because the risen Christ dies no more (Rom~6:9), His living Body cannot be separated from His Blood, human soul, or divinity. By \textbf{natural and hypostatic concomitance} (\latin{ex naturali concomitantia}), wherever Christ's Body is present under bread, His Blood, soul, and Godhead are wholly present with it; and under wine, His Body, soul, and Godhead are equally wholly present (Trent, Sess.~XIII, ch.~3; \emph{ST}~III, q.~76, aa.~1--2; CCC~1390). The whole Christ is present whole and entire in every intact particle and drop.

\paragraph{5. Sacramental mode of presence and abiding duration.}
Christ is present ``after the manner of substance'' (\latin{per modum substantiae}), and ``in no way \dots\ locally in this sacrament'' (\emph{ST}~III, q.~76, a.~5). His glorified dimensions are not spatially crammed inside the Host, nor does He leave heaven or multiply when consecrated on thousands of altars. Because transubstantiation alters the substance of the gifts themselves, Christ's Real Presence endures as long as the consecrated species remain intact (Trent, Sess.~XIII, ch.~5, can.~4; CCC~1377), ceasing only when the species are digested or corrupted beyond the nature of bread and wine. This abiding presence grounds reservation in the tabernacle, Viaticum for the dying, and supreme adoration (\latin{latria}) of the Blessed Sacrament.

\begin{threestable}{Metaphysical Category}{Exact Catholic Determination}{Error Excluded}
Transubstantiation & Unique substantial conversion of the entire substance of bread and wine into Christ's Body and Blood by divine power (Lateran~IV; Trent~XIII, ch.~4; \emph{ST}~III, q.~75, a.~4) & Consubstantiation (bread remaining), annihilation without conversion, or mere symbolic figure\\
Words of Consecration & Instrumental sacramental form spoken \latin{in persona Christi} within the Eucharistic Prayer (\emph{ST}~III, q.~78, aa.~1, 4) & Mere verbal re-labeling (``transignification'' or ``transfinalization'' apart from ontological change; \emph{Mysterium Fidei})\\
Sensible Accidents (\latin{Species}) & Color, taste, texture, weight, and dimensive quantity remain unaltered, sustained miraculously without their former substance (\emph{ST}~III, q.~77, aa.~1--2) & Physical alteration of appearances, sensory illusion, or bodily division of Christ\\
\latin{Ex vi verborum} vs.\ Concomitance & Body under bread and Blood under wine \latin{ex vi verborum}; whole living Christ wholly present under each intact species by concomitance (\emph{ST}~III, q.~76, aa.~1--2; CCC~1390) & Dividing the risen Christ, or holding that Communion under one kind gives less than the whole Christ\\
Mode and Duration & Present ``after the manner of substance,'' not locally circumscribed (\emph{ST}~III, q.~76, a.~5); abides while species remain intact (Trent~XIII, can.~4; CCC~1377) & Spatial localization, multiplication of Christ's body, or restricting presence solely to the moment of eating\\
\end{threestable}

\subsection{The Sacrificial Identity of the Mass and Calvary}
\label{subsec:eucharist-sacrifice}

The Eucharist is inseparably both sacrament and sacrifice: ``This sacrament is both a sacrifice and a sacrament'' (\emph{ST}~III, q.~79, a.~5). The Council of Trent in Session~XXII (1562, chs.~1--4, cans.~1--9), Vatican~II (\emph{Sacrosanctum Concilium}~47--48), St.~John Paul~II (\emph{Ecclesia de Eucharistia}~11--20), Pope Benedict~XVI (\emph{Sacramentum Caritatis}~10--11), and CCC~1362--1372, 1382 anchor this doctrine in the all-sufficient Cross.

Because Christ offered Himself ``once'' for eternal redemption (Heb~7:27), the Mass never repeats Calvary or adds a new atonement. Instituted at the Last Supper, it is the visible, unbloody sacrifice that sacramentally re-presents (\latin{anamnesis}) the bloody sacrifice of the Cross, perpetuates its memory, and applies its saving power for the remission of daily sins (Trent, Sess.~XXII, ch.~1). Trent defines their strict identity (Sess.~XXII, ch.~2): the \textbf{Victim} is one and the same Jesus Christ; the \textbf{Principal Offerer} is one and the same Christ now offering through the ministry of ordained priests; and only the \textbf{manner of offering} differs---bloody once on Calvary, unbloody on the altar through the separate consecration of bread and wine signifying the separation of Body and Blood in death (\emph{ST}~III, q.~83, a.~1).

The Mass is truly \textbf{propitiatory}, offered to God alone---never to the saints, ``but to God alone, who crowned them'' (\latin{non tamen illis sacrificium offerri docet, sed Deo soli, qui illos coronavit}; Trent, Sess.~XXII, ch.~3)---and avails for both the living and the faithful departed (Trent, Sess.~XXII, ch.~2, can.~3; Cyril, \emph{Myst.~Cat.}~V.8, 10). Into Christ's offering the whole Church joins her prayers and life (\emph{De civitate Dei}~X.6; \emph{Sacrosanctum Concilium}~48). And because the altar of sacrifice is equally the Lord's table, sacrificial memorial and paschal banquet are inseparable: the Victim offered to the Father is given back to the faithful as heavenly food (CCC~1382).

\subsection{Matter, Form, Anaphora, Minister, Subject, Effects, and Ends}
\label{subsec:eucharist-constitution}

\paragraph{Matter.}
Christ established the remote matter at the Last Supper: \textbf{wheaten bread} and \textbf{natural wine of the grape} (\emph{ST}~III, q.~74, aa.~1, 3, 5; CIC~1983, c.~924; CCEO~1990, c.~706). Non-wheat grains or liquids other than genuine grape wine (outside lawfully authorized mustum for priests unable to take alcohol) are invalid. By apostolic precept for liceity, a drop of water is mingled with the wine, signifying the blood and water from Christ's side and the union of the faithful with Christ (\emph{ST}~III, q.~74, a.~6; Trent, Sess.~XXII, can.~9). The Council of Florence (\emph{Laetentur caeli}, 1439, DS~1303) affirmed that Christ's Body is truly consecrated in both unleavened and leavened wheaten bread, each Church keeping its ancient rite: unleavened bread (\latin{azymes}) in the Latin, Armenian, and Maronite Churches (CIC~1983, c.~926), and leavened bread (\latin{artos}) in the Byzantine and other Eastern Churches (CCEO~1990, c.~707~\S1).

\paragraph{Form and the Unity of the Anaphora (East and West).}
In the Latin scholastic and liturgical tradition, Christ's words of institution pronounced \latin{in persona Christi} (``in the person of Christ'') within the Eucharistic Prayer---in the Latin \emph{editio typica} of the Roman Missal, \latin{Hoc est enim Corpus meum, quod pro vobis tradetur} over the bread and \latin{Hic est enim calix Sanguinis mei novi et aeterni testamenti, qui pro vobis et pro multis effundetur in remissionem peccatorum} over the chalice, identifying the consecrated gifts as Christ's Body given for us and His Blood of the new and eternal covenant shed for us and for many unto the remission of sins---constitute the essential consecratory form (\emph{ST}~III, q.~78, aa.~1--4; cf.~CDWDS, 17~October 2006 circular letter on rendering \latin{pro multis}, ``for many''). Speaking in the first person, the priest lends his voice directly to Christ Himself.

In the Eastern Catholic Churches---praying the ancient Byzantine Anaphoras of St.~John Chrysostom and St.~Basil, the West Syriac Anaphora of St.~James, the East Syriac Anaphora of Addai and Mari, and the Alexandrian and Armenian anaphoras---consecration is contemplated within the unified movement of the entire \textbf{Anaphora}. Here Christ's institution words and the \textbf{epiclesis} invoking the Holy Spirit upon the gifts and the assembly (seen in St.~Cyril, \emph{Myst.~Cat.}~V.7, and St.~John of Damascus, \emph{De fide orth.}~IV.13) form one indivisible Trinitarian act. Both Latin Christological focus and Eastern pneumatological breadth confess the same mystery: the Father sends the Spirit to fulfill the creative word of His Son.

\paragraph{Minister and Subject.}
Valid consecration belongs exclusively to a validly ordained \textbf{bishop or priest} acting \latin{in persona Christi Capitis} (``in the person of Christ the Head''; Lateran~IV, \emph{Firmiter}, DS~802; Trent, Sess.~XXII, can.~2; \emph{ST}~III, q.~82, aa.~1, 5; CIC~1983, c.~900~\S1; CCEO~1990, c.~699~\S1). No deacon or layperson can ever consecrate. Ordinary ministers of distribution are the bishop, priest, and deacon (CIC~1983, c.~910~\S1; CCEO~1990, c.~709~\S1), with instituted acolytes or mandated extraordinary ministers assisting in genuine Latin pastoral need (CIC~1983, c.~910~\S2).

Any living, validly baptized person not barred by law is a capable subject (CIC~1983, c.~912). A believer who has reached reason must approach in the \textbf{state of sanctifying grace}: receiving while conscious of unrepented mortal sin contradicts the sacrament of charity and commits sacrilege (1~Cor~11:27--29; \emph{ST}~III, q.~80, aa.~1, 4; Trent, Sess.~XIII, ch.~7; CIC~1983, c.~916; CCEO~1990, c.~711), requiring prior sacramental Confession. Reverent preparation also includes the Eucharistic fast (one hour in Latin law, CIC~1983, c.~919; governed by particular law in the East, CCEO~1990, c.~713). Latin law admits children at the age of discretion with proper preparation (CIC~1983, cc.~913--914), whereas Eastern common law directs admission of baptized and chrismated infants according to each Church's liturgical books and particular law (CCEO~1990, cc.~697, 710). Because of real concomitance, the whole Christ is received under either kind alone, while Communion under both kinds expresses the sign of the banquet more fully when permitted by liturgical discipline (Trent, Sess.~XXI, chs.~1--3; CIC~1983, c.~925; CCC~1390).

\begin{fourstable}{Dimension}{Common Catholic Doctrine}{Latin Church (CIC~1983)}{Eastern Churches (CCEO~1990)}
Matter & Wheaten bread and natural grape wine with a drop of water (\emph{ST}~III, q.~74; Florence, \emph{Laetentur caeli}) & Unleavened wheaten bread (\latin{azymes}) and grape wine (cc.~924, 926) & Leavened bread (\latin{artos}, Byzantine/most East) or unleavened (Armenian/Maronite) and grape wine (cc.~706--707)\\
Consecration & Christ acts through ordained priest in Trinitarian Eucharistic Prayer & Christ's institution words spoken \latin{in persona Christi} highlighted as form within Eucharistic Prayer & Unified Anaphora joining institution narrative and Spirit-epiclesis (Cyril, \emph{Myst.~Cat.}~V.7; Damascus, \emph{De fide orth.}~IV.13)\\
Minister & Consecration reserved to validly ordained bishop or priest alone (Lateran~IV; Trent~XXII) & Consecrator: bishop or priest (c.~900). Distribution: bishop, priest, deacon (c.~910~\S1); extraordinary ministers in need (c.~910~\S2) & Consecrator: bishop or priest (c.~699~\S1). Distribution: priest, or deacon per particular law (c.~709~\S1)\\
Subject and Communion & Validly baptized; state of grace required for responsible adults (1~Cor~11:27--29); whole Christ under each kind & First Eucharist at age of discretion (cc.~913--914); one-hour fast (c.~919); one or both kinds (c.~925) & Infant Communion with B--C--E per liturgical books (cc.~697, 710); Communion normally under both kinds by intinction/spoon\\
\end{fourstable}

\paragraph{Sacramental Structure, Effects, and Order of Ends.}
The threefold architecture of the Eucharist stands at the apex of all seven mysteries (\emph{ST}~III, q.~73, aa.~3--4; q.~79, aa.~1--6):
\begin{itemize}
  \item \textbf{\latin{Sacramentum tantum}} (the outward sign alone): the consecrated sensible species of bread and wine.
  \item \textbf{\latin{Res et sacramentum}} (the intermediate reality and sign): the \textbf{true Body and Blood of Jesus Christ} substantially present under those species. The Eucharist imprints no character and is repeatable daily; its \latin{res et sacramentum} is not a created spiritual potency, but the living Lord Himself enduring while the species remain intact.
  \item \textbf{\latin{Res tantum}} (the ultimate proper grace): \textbf{the unity of the Mystical Body in charity} (\emph{ST}~III, q.~73, a.~3; Augustine, \emph{In Ioann. Evang.}~Tr.~26.13)---intimate abiding union with Christ (John~6:57 DR [6:56 Gk]), increase of sanctifying grace and love, remission of venial sins, preservation from future mortal sins, care for the poor, and the pledge of bodily resurrection and eternal glory (\emph{ST}~III, q.~79, aa.~1--6; Benedict~XVI, \emph{Sacramentum Caritatis}~70--71; CCC~1391--1405).
\end{itemize}
Across time, St.~Thomas maps its three venerable names to past, present, and future (\emph{ST}~III, q.~73, a.~4): as memorial of Christ's Passion it is \textbf{Sacrifice}; as causing present ecclesial unity it is \textbf{Communion} or \latin{Synaxis}; and as foretaste of eternal glory it is \textbf{Viaticum} (``food for the journey'') and \textbf{Eucharist} (``good grace'' or thanksgiving).

Its ends stand in clear hierarchical order:
\begin{enumerate}
  \item \textbf{Primary proper end:} spiritual nourishment of the baptized soul through intimate union with Christ and sacrificial communion in His Paschal Mystery.
  \item \textbf{Intrinsic secondary ends:} building up the unity of the Mystical Body, rendering supreme worship and thanksgiving to the Father, offering propitiation for the living and the dead, and healing venial sin while fortifying against mortal sin.
  \item \textbf{Contingent ends:} spiritual consolation in trial, apostolic zeal in daily life, and the special strengthening of the dying through Viaticum.
  \item \textbf{Ultimate end:} eternal life, bodily resurrection on the last day (John~6:55 DR [6:54 Gk]), and everlasting fellowship at ``the marriage supper of the Lamb'' (Apoc~19:9; CCC~1402--1405).
\end{enumerate}

All the streams of Christian Initiation empty into this ocean. Washed in the waters of Baptism and sealed with the fragrant chrism of the Spirit, the Christian steps to the altar, answers ``Amen'' to the true Flesh and Blood of the Word made flesh, and receives already under humble veils the eternal life for which he was created.
